\appendix

\section{Negative mechanism results}
\label{app:negative}

This appendix first records the prototype treatment conditions and injection magnitudes, then gives class reliability and treatment details, and finally documents the isolation diagnostics.

This appendix holds the measurements behind Section~\ref{sec:res-prototype}. They are
reported in full because the case for changing the feedback model's consultant rests on them, and
because each one closes a line of investigation that would otherwise look open.

\FloatBarrier
\subsection{Conditions A, B and C}
\label{app:abc}

Three configurations of the feedback model were run at three seeds on both datasets, with the fold
partition fixed. Condition A is the canonical prototype-consultant feedback model. Condition B trains
the selector head whose entropy ranks edges and calibrates the consultant's temperature per
fold instead of pinning it. Condition C re-selects the two knobs on top of B. The GNN model and fusion model rungs are unaffected by any of the three and are shown for reference in
Table~\ref{tab:abc}.

\begin{table}[htbp]
\centering
\caption[Pooled out-of-fold macro-F1, mean over three seeds, by condition]{Pooled out-of-fold macro-F1, mean over three seeds, by condition.}
\label{tab:abc}
\small
\begin{tabular}{llccc}
\toprule
Dataset & Rung & A (canonical) & B (signals fixed) & C (B, re-selected) \\
\midrule
\multirow{3}{*}{UNSW}
 & feedback model          & 0.7644 & 0.7610 & 0.7809 \\
 & fusion model          & 0.7793 & 0.7793 & 0.7793 \\
 & GNN model & 0.7437 & 0.7437 & 0.7437 \\
\midrule
\multirow{3}{*}{ToN}
 & feedback model          & 0.4521 & 0.4147 & 0.4630 \\
 & fusion model          & 0.4102 & 0.4102 & 0.4102 \\
 & GNN model & 0.4336 & 0.4336 & 0.4336 \\
\bottomrule
\end{tabular}
\end{table}

The decision rule was fixed before the runs. It returns the same verdict in all three
conditions: the feedback model is not separated from the fusion model on either dataset. On NF-UNSW-NB15 the
two-level interval for the feedback model minus the fusion model is $-0.0145$ [$-0.0556$, $+0.0276$] in A, $-0.0181$
[$-0.0596$, $+0.0242$] in B and $+0.0018$ [$-0.0406$, $+0.0444$] in C. On NF-ToN-IoT it is
$+0.0433$ [$-0.0376$, $+0.1199$], $+0.0065$ [$-0.0782$, $+0.0788$] and $+0.0531$ [$-0.0360$,
$+0.1340$]. Condition C carries the highest point estimate on both datasets. That is not a
separation, and its selection curves are flat, with NF-UNSW-NB15's argmax landing on the
boundary of the swept range at $k = 35$.

Both defects were genuinely removed in B. That is what makes the negative result
informative rather than a failed implementation. On NF-UNSW-NB15 the selector head's macro-F1
moved from the range 0.019 to 0.036 up to the range 0.646 to 0.684, mean disagreement fell
from 0.8975 to between 0.524 and 0.548, the consultant's maximum softmax probability rose
from 0.102 to 0.487, and the temperature fell from about 9.4 to 0.105. On NF-ToN-IoT the
selector head moved from 0.029 to 0.065 up to 0.293 to 0.333, disagreement fell from 0.8988
to between 0.580 and 0.638, maximum probability rose from 0.101 to 0.500, and temperature
fell from about 9.5 to 0.030. A better-calibrated consultant and a genuinely
uncertainty-ranked selector made the feedback model worse on both datasets.

\FloatBarrier
\subsection{What the canonical feedback model was actually injecting}
\label{app:magnitudes}

Table~\ref{tab:magnitudes} gives the mean absolute magnitude added to the 39-wide encoded
edge vector on flagged edges, per seed. The encoded features have unit variance, so these numbers are
directly interpretable as a fraction of a standard deviation.

\begin{table}[htbp]
\centering
\caption[Injected bias magnitude on flagged edges, per seed, by condition]{Injected bias magnitude on flagged edges, per seed, by condition.}
\label{tab:magnitudes}
\small
\begin{tabular}{lll}
\toprule
Condition & NF-UNSW-NB15 & NF-ToN-IoT \\
\midrule
A (canonical)        & 0.077 / 0.060 / 0.061 & 0.023 / 0.021 / 0.023 \\
B (signals fixed)    & 1.228 / 1.171 / 1.218 & 2.796 / 2.854 / 2.762 \\
C (B, re-selected)   & 1.103 / 1.100 / 1.158 & 2.845 / 2.893 / 2.992 \\
\bottomrule
\end{tabular}
\end{table}

Under the canonical prototype configuration the mechanism was injecting between 0.02 and 0.08
of a standard deviation. Calibrating the temperature raised that by a factor of roughly 17 on
NF-UNSW-NB15 and roughly 125 on NF-ToN-IoT. Two things follow. The prototype feedback model's advantage
over feedback-off was late fusion rather than feedback. The injection scale is also close to
redundant with the learned bias strength, since condition C moved the scale from 2.0 to 5.0 and
from 20.0 to 10.0 while the end-to-end magnitude barely moved.

\FloatBarrier
\subsection{Per-class reliability of the embedding-only consultant}
\label{app:reliability}

The reliability-weighting treatment computes, for each class, the precision of the prototype
consultant's argmax on the fold's training edges. The weights are a more useful result than
the macro-F1 deltas they produced, because they quantify the consultant's weakness per class
from training labels alone. Table~\ref{tab:reliability} gives them.

\begin{table}[htbp]
\centering
\caption[Prototype precision on training edges]{Mean per-class precision of the prototype consultant on training edges, over five
folds. Each dataset's classes are listed in its own label-mapping order.}
\label{tab:reliability}
\small
\begin{tabular}{clclc}
\toprule
 & \multicolumn{2}{c}{NF-UNSW-NB15} & \multicolumn{2}{c}{NF-ToN-IoT} \\
\cmidrule(lr){2-3}\cmidrule(lr){4-5}
Index & class & precision & class & precision \\
\midrule
0 & Normal         & 1.000 & Benign     & 1.000 \\
1 & Analysis       & 0.852 & backdoor   & 0.071 \\
2 & Backdoors      & 0.976 & ddos       & 0.179 \\
3 & DoS            & 0.893 & dos        & 0.021 \\
4 & Exploits       & 0.961 & injection  & 0.570 \\
5 & Fuzzers        & 0.662 & mitm       & 0.740 \\
6 & Generic        & 0.895 & password   & 0.226 \\
7 & Reconnaissance & 0.984 & ransomware & 0.020 \\
8 & Shellcode      & 1.000 & scanning   & 0.095 \\
9 & Worms          & 0.960 & xss        & 0.089 \\
\bottomrule
\end{tabular}
\end{table}

On NF-UNSW-NB15 six of ten classes sit at or above 0.9 and the minimum is 0.662, so
reweighting has little to reorder. On NF-ToN-IoT six of ten sit below 0.226: \texttt{ransomware}
at 0.020, \texttt{dos} at 0.021, \texttt{backdoor} at 0.071, \texttt{xss} at 0.089,
\texttt{scanning} at 0.095 and \texttt{ddos} at 0.179, with \texttt{password} a seventh sitting
at 0.226 itself. That is the same weakness the isolation experiment measures as the binding
constraint, now quantified per class.

\FloatBarrier
\subsection{The five treatments}
\label{app:treatments}

% Qualitative summary, written by hand. Not generated from a contract.
\begin{table}[htbp]
\centering
\caption{Five treatments applied to the embedding-only consultant. None was kept. Rows marked
seed 42 carry no interval across seeds and are single-seed orderings only.}
\label{tab:treatments}
\footnotesize
\begin{tabularx}{\textwidth}{@{}l X X c@{}}
\toprule
Treatment & What it changed & Result & Kept \\
\midrule
Selector-head supervision & the ranking head is trained as a classifier & 3 seeds: loop worse on both datasets & no \\
Consultant temperature & per-fold calibration instead of a pinned $T$ & 3 seeds: loop worse, jointly with the above & no \\
Reliability weighting & gate and magnitude scaled by per-class train precision & seed 42: $+0.0057$ / $+0.0094$, below the 0.02 rule & no \\
Advice format & one-hot and softmax at the oracle's $\pm 4$ nats & seed 42: every arm below baseline on both & no \\
Cross-fitted advice & train rows replaced by inner out-of-fold logits & seed 42: $-0.0246$ / $-0.0338$ against in-fold & no \\
\bottomrule
\end{tabularx}
\end{table}


The advice-format arm produced a counter-intuitive measurement worth stating on its own. Handing
the consultant the oracle's saturated format made the input magnitude larger and the injected
magnitude smaller, from 0.0767 to 0.0183 on NF-UNSW-NB15 and from 0.0234 to 0.0202 on
NF-ToN-IoT. The projection is learned, so it adapts to a larger and lower-variance input by
injecting less. Whatever the oracle's advantage is, it is not the shape of the vector it presents.

The cross-fitting arm tested the feedback model's trust directly. If in-fold optimism in the advice were
what the feedback model was exploiting, honest advice should lower the scalar the feedback model learns for how hard
to push. Across four settings the learned trust scalar moved by at most 0.0205, from 1.5722 to
1.5664 and 1.5227 to 1.5022 on NF-UNSW-NB15 and from 1.4768 to 1.4652 and 1.4842 to 1.4899 on
NF-ToN-IoT, moving the wrong way in the last of those. What did move was the injected magnitude,
from 1.3430 to 1.1752 and from 1.3291 to 1.0745. The model absorbs the change in the projection
rather than in the parameter that is supposed to represent trust.

\FloatBarrier
\subsection{Isolation diagnostics and prediction changes}
\label{app:isolation-detail}

The isolation measurements use the earlier consultation percentiles 31 on NF-UNSW-NB15 and 25 on NF-ToN-IoT, with output fusion disabled. They should not be read as a direct ablation of the final percentiles 29 and 16. The paired intervals resample (seed, fold) pairs and differ from the two-level edge bootstrap used for the model ladder.

The oracle headroom, computed from the plain three-seed means, is $+0.0345$ on NF-UNSW-NB15 and $+0.1090$ on NF-ToN-IoT. The captured shares in Table~\ref{tab:mechanism} use these values, rather than the paired-bootstrap mean differences. The prototype captures $-4.5\%$ and $11.9\%$, neither separated from zero; the trained head captures $12.5\%$ and $-38.0\%$, with the latter a separated loss.

The NF-UNSW-NB15 headroom comes with a caveat. At seed 42 alone the oracle arm scores
0.7522 against a control of 0.7543, a headroom of $-0.0021$, so a captured share is not
computable there. Because the $+0.0345$ is a mean over three seeds that do not agree in sign,
single-seed shares on NF-UNSW-NB15 are not reported.

A cross-fitted trained-head diagnostic at seed 42 moves the difference against control from $-0.0230$ to $-0.0654$ on NF-UNSW-NB15 and from $-0.0165$ to $+0.0070$ on NF-ToN-IoT. The single-seed result does not establish a reliable improvement.

Churn, the fraction of edges whose predicted class changes
between consecutive passes, measured in the canonical feedback model over every edge in the graph,
training edges included, is 0.1372 on NF-UNSW-NB15 and 0.1246 on NF-ToN-IoT as three-seed means,
with per-seed values from 0.0938 to 0.1675 and from 0.1070 to 0.1402. The injection path is
alive and moves predictions on every pass. This measures prediction changes, not accuracy gains. The churn figures are in
\texttt{results/feedback\_churn.json}, and Section~\ref{sec:res-attention} contrasts them with the
attention variant, whose churn is exactly zero.

The churn summary is saved in \texttt{results/feedback\_churn.json}, but its source traces are in the local, untracked \texttt{tmp/multiseed\_v2} directory. The summary is available; independent recalculation of these particular values requires those traces.

\clearpage
\section{Supporting tables}
\label{app:tables}

This appendix first provides the per-class and selection tables, then examines consultant complementarity, and finally gives detailed model comparisons and baseline qualifications.

\FloatBarrier
\subsection{Per-class F1}
\label{app:perclass}

% Generated by docs/research_report/tables/make_tables.py. Do not edit by hand.
\begin{table}[htbp]
\centering
\caption[Per-class F1 and edge counts]{Per-class F1, mean over three seeds, with edge counts. The feedback model is the canonical one, consulting the trained head. NF-ToN-IoT's \texttt{dos} and \texttt{ransomware} are excluded from the metric and omitted here, and F1 is computed on the rows the metric scores, so each column's mean is that rung's macro-F1 in Table~\ref{tab:ladder}.}
\label{tab:perclass}
\small
\begin{tabular}{lrccc}
\toprule
Class & edges & GNN & fusion & feedback \\
\midrule
\multicolumn{5}{l}{\emph{NF-UNSW-NB15}} \\
Normal                   &  311 & 0.9129 & 0.9152 & 0.9543 \\
Analysis                 &   25 & 0.4327 & 0.5886 & 0.6196 \\
Backdoors                &   40 & 0.5929 & 0.6578 & 0.8019 \\
DoS                      &   40 & 0.5831 & 0.6446 & 0.7424 \\
Exploits                 &   40 & 0.9309 & 0.9583 & 0.9580 \\
Fuzzers                  &   40 & 0.6207 & 0.7330 & 0.7749 \\
Generic                  &   40 & 0.6591 & 0.5409 & 0.6899 \\
Reconnaissance           &   40 & 0.9110 & 0.8916 & 0.8731 \\
Shellcode                &   40 & 0.8669 & 0.9428 & 0.9959 \\
Worms                    &   40 & 0.9272 & 0.9205 & 0.9311 \\
\midrule
\multicolumn{5}{l}{\emph{NF-ToN-IoT}} \\
Benign                   & 1746 & 0.9599 & 0.9280 & 0.9835 \\
\texttt{backdoor}        &   16 & 0.3817 & 0.2439 & 0.4720 \\
\texttt{ddos}            &   35 & 0.1901 & 0.2498 & 0.3000 \\
\texttt{injection}       &  116 & 0.5436 & 0.4851 & 0.6277 \\
\texttt{mitm}            &  157 & 0.8838 & 0.7082 & 0.8566 \\
\texttt{password}        &   25 & 0.0401 & 0.2996 & 0.2134 \\
\texttt{scanning}        &   13 & 0.0969 & 0.2829 & 0.3211 \\
\texttt{xss}             &   12 & 0.3728 & 0.0844 & 0.2612 \\
\bottomrule
\end{tabular}
\end{table}


\FloatBarrier
\subsection{Knob selection curves}
\label{app:knobs}

Table~\ref{tab:knobcurves} summarises the validation-only searches. The ranges show how little the score changes across the tested settings and identify the boundary-selected scale.

% Generated by docs/research_report/tables/make_tables.py. Do not edit by hand.
\begin{table}[htbp]
\centering
\caption{Knob selection curves for the trained-head consultant, mean validation macro-F1 over selection seeds 42, 1 and 2. Every span is smaller than the rungs' own seed-to-seed spread, so the argmax is not distinguishable from its neighbours.}
\label{tab:knobcurves}
\footnotesize
\resizebox{\textwidth}{!}{%
\begin{tabular}{llccccccc}
\toprule
Dataset & Knob & candidates & range & validation macro-F1 & span & selected & on boundary \\
\midrule
NF-UNSW-NB15 & $k$ (entropy percentile) & 21 & 15--35 & 0.8238--0.8277 & 0.0038 & 29 & no \\
NF-UNSW-NB15 & injection scale & 6 & 0.5--20 & 0.8252--0.8277 & 0.0025 & 2 & no \\
NF-ToN-IoT & $k$ (entropy percentile) & 21 & 15--35 & 0.5170--0.5232 & 0.0062 & 16 & no \\
NF-ToN-IoT & injection scale & 6 & 0.5--20 & 0.5196--0.5232 & 0.0037 & 20 & yes \\
\bottomrule
\end{tabular}
}
\end{table}


\FloatBarrier
\subsection{Consultant complementarity on the flagged set}
\label{app:complementarity}

The feedback model consults the edges the GNN model is least certain about, so what matters is not a
consultant's accuracy over the whole graph but its accuracy where it is actually asked.
Table~\ref{tab:complementarity} gives that, at seed 42.

\begin{table}[htbp]
\centering
\caption[Consultant accuracy on uncertain edges]{Accuracy on the entropy-flagged edge set at seed 42, for the GNN model, each
consultant, and the resulting feedback model.}
\label{tab:complementarity}
\small
\begin{tabular}{llccc}
\toprule
Dataset & Consultant & GNN model & consultant & resulting feedback model \\
\midrule
\multirow{2}{*}{NF-UNSW-NB15}
 & prototype    & 0.547 & 0.647 & 0.658 \\
 & trained head & 0.547 & 0.795 & 0.779 \\
\midrule
\multirow{2}{*}{NF-ToN-IoT}
 & prototype    & 0.343 & 0.284 & 0.361 \\
 & trained head & 0.343 & 0.648 & 0.628 \\
\bottomrule
\end{tabular}
\end{table}

The NF-ToN-IoT prototype row is the clearest statement of the problem. On the edges the feedback model
selects, that consultant is less accurate than the model it is supposed to correct.

Table~\ref{tab:disagreement} gives what happens where the two disagree, the only
place a consultant can change an answer.

\begin{table}[htbp]
\centering
\caption[Consultant correctness on disagreements]{Among entropy-flagged edges where the consultant and the GNN model disagree, how
often the consultant is right. Seed 42, at the canonical entropy percentile for each dataset.}
\label{tab:disagreement}
\small
\begin{tabular}{llcccc}
\toprule
 & & \multicolumn{2}{c}{flagged set} & \multicolumn{2}{c}{confidence-gated half} \\
\cmidrule(lr){3-4}\cmidrule(lr){5-6}
Dataset & Consultant & right & wrong & right & wrong \\
\midrule
\multirow{2}{*}{NF-UNSW-NB15}
 & prototype    & 38 & 45 & 13 & 5 \\
 & trained head & 55 & 22 & 17 & 4 \\
\midrule
\multirow{2}{*}{NF-ToN-IoT}
 & prototype    & 45 & 178 & 28 & 63 \\
 & trained head & 129 & 75 & 79 & 13 \\
\bottomrule
\end{tabular}
\end{table}

The gate is doing its job for the trained head and not for the prototype. Restricting to the
consultant's own most confident half lifts the trained head from 55 against 22 to 17 against 4
on NF-UNSW-NB15 and from 129 against 75 to 79 against 13 on NF-ToN-IoT. It does not rescue the
prototype on NF-ToN-IoT, where confident disagreements are wrong more than twice as often as
they are right. A gate can only concentrate a consultant's competence; it cannot manufacture it.

Both tables are recomputed from the committed pooled out-of-fold tensors, by a script that
writes its own artifact:

\begin{quote}
\texttt{scripts/close\_complementarity\_split.py} \\
\texttt{results/consultant\_complementarity.json}
\end{quote}

\FloatBarrier
\subsection{Model comparisons and metric interpretation}
\label{app:ladder-detail}

Macro-F1 weights each evaluated class equally. Accuracy and weighted F1 give more influence to frequent classes, so on NF-ToN-IoT they are strongly influenced by benign traffic, which accounts for 82\% of scored edges. The component comparison in Table~\ref{tab:metrics} reports all three metrics, including the head alone. A high accuracy can therefore coexist with weak rare-class performance.

The feedback model's position depends on its consultant. With the prototype scorer it is above the GNN model and the semantic model on both datasets, below the fusion model on NF-UNSW-NB15 by 0.0149, and above
the fusion model on NF-ToN-IoT by 0.0419. With the trained head it is the highest of the four stages on both, but the head alone has a higher mean and its difference from feedback is not separated. The change
is $+0.0697$ on NF-UNSW-NB15 and $+0.0523$ on NF-ToN-IoT. Two things changed together and the
report does not separate them: the consultant, and the selected entropy percentile, which moved
from 31 to 29 and from 25 to 16.

Both configurations stored their pooled out-of-fold predictions at all three seeds, so the
change carries an interval computed the same way as every other difference here. It is
$+0.0701$ [$+0.0353$, $+0.1057$] on NF-UNSW-NB15, separated, and $+0.0511$ [$-0.0082$,
$+0.1110$] on NF-ToN-IoT, not separated. The seed-matched interval on NF-ToN-IoT does exclude
zero, at $+0.0509$ [$+0.0022$, $+0.1013$]. That is the same pattern the fusion model against the GNN model shows below: a difference that survives resampling edges and stops surviving once
training-seed variance is admitted. The
interval measures the change as a whole. It does not attribute it to the consultant, because
the entropy percentile moved with it.

The trained consultant was also scored standing alone. The
trained head reaches $0.8374 \pm 0.0048$ on NF-UNSW-NB15 and $0.5081 \pm 0.0073$ on NF-ToN-IoT.
It is separated above the GNN model on both, at $+0.0946$ [$+0.0403$, $+0.1470$] and $+0.0725$
[$+0.0127$, $+0.1313$]. The feedback model is \emph{not} separated from it: $-0.0032$
[$-0.0228$, $+0.0128$] on NF-UNSW-NB15 and $-0.0032$ [$-0.0342$, $+0.0299$] on NF-ToN-IoT, the
sign is unstable across the three seeds, and the head carries the higher mean on both. We report
this because the consultant experiment is the central manipulation of this study and a reader
cannot weigh it without knowing where the consultant sits on its own. It is shown as a component in Table~\ref{tab:metrics}, rather than as a parallel stage of the ladder, and Section~\ref{sec:disc-contribution} reads it against the isolation experiment.

Figure~\ref{fig:comparisons} and Table~\ref{tab:comparisons} report the stage comparisons for both consultants. Under the
prototype scorer, no comparison between the feedback model and a structural rung is separated on
either dataset. Every one of those intervals crosses zero, including the two where the feedback
model carries the higher mean. The one prototype-consultant difference that is separated is the
feedback model against the semantic model on NF-ToN-IoT, at $+0.1727$ [$+0.1326$, $+0.2174$],
where the semantic model scores 0.2785 against the GNN model's 0.4336 and clearing it says
nothing about the mechanism. On NF-UNSW-NB15 even that comparison crosses zero, at $+0.0287$
[$-0.0061$, $+0.0631$]. In that configuration the answer to RQ3 is that the ladder is an
ordering of means and no more.

One result is negative regardless of consultant. The fusion model is not separated from the GNN model:
positive on NF-UNSW-NB15 and negative on NF-ToN-IoT, with an interval crossing zero either way.
The seed-matched interval for NF-UNSW-NB15 does exclude zero, $+0.0356$ [$+0.0038$, $+0.0678$].
That is the clearest illustration in this report of what the two-level procedure is for: once
seed variance enters, the same difference stops being separated. Under the pre-registered rule,
the fusion model versus the GNN model is not separated, and the answer to RQ2 is negative on both
datasets.

\FloatBarrier
\subsection{Baseline scoring asymmetry and comparison details}
\label{app:baseline-detail}

The feedback comparisons below remain subject to the component check: the trained head alone has a higher mean macro-F1 on both datasets, with no separated difference from feedback (Table~\ref{tab:metrics}).

The collapse of E-GraphSAGE on NF-UNSW-NB15 is a representation effect with a testable
prediction attached, not evidence that the architecture is weak. Its edge classifier reads only the
concatenation $[h_u \Vert h_v]$ of the two endpoint embeddings. An edge's own attributes never
reach it, so parallel edges between one address pair cannot be told apart. That predicts a
collapse
proportional to how common parallel edges are. On NF-UNSW-NB15, 385 of 656 edges are
indistinguishable by endpoints, that is 58.7\%, and the architecture scores 0.1194. On
NF-ToN-IoT only 167 of 2{,}127 edges are, that is 7.9\%, and the same architecture scores
0.4141. The prediction and the inversion agree. Figure~\ref{fig:ceiling} shows the mechanism
beside the confirmation.
The ceiling is also not a training-schedule artifact: trained to the authors' full 4{,}999
epochs on fold 0, validation macro-F1 plateaus near 0.11 and peaks at 0.1384 at epoch 1{,}400.

The lower rungs are weaker still. On NF-ToN-IoT neither the GNN model nor the fusion model is separated
from E-GraphSAGE as published, at $+0.0203$ [$-0.0230$, $+0.0633$] and $-0.0061$ [$-0.0797$,
$+0.0777$]. On NF-UNSW-NB15 neither the GNN model nor the semantic model is separated from
TE-G-SAGE with our node features, at $+0.0249$ and $+0.0163$.

The NF-ToN-IoT baselines were scored under a looser rule than the rungs they are compared
against, and the margin has to be read against that. As Section~\ref{sec:method-protocol}
states, these models and the fusion model could predict \texttt{dos} or \texttt{ransomware} for
an edge the metric scores, losing it, while the GNN, semantic and feedback models could not.
Across the three seeds that happened on 3, 374 and 6 edges for the fusion model, on 58, 63 and
47 for E-GraphSAGE as published, on 66, 28 and 60 for E-GraphSAGE with our node features, and on
between 29 and 206 for the three TE-G-SAGE configurations. Replacing every one of those
predictions with the true label gives an upper bound on what the mask could have been worth, and
that bound is more generous than any mask could be, because no mask redirects an edge to the
right class. The bound stays below the feedback model at every seed for the fusion model and for
five of the six baselines. The exception is E-GraphSAGE with our node features at seed 1, where
it reaches 0.5206 against the feedback model's 0.5012. Table~\ref{tab:vsbaselines} already
reports that comparison as not separated, at $+0.0667$ [$-0.0181$, $+0.1505$], so no claim in
this report depends on the asymmetry. Every count and both macro-F1 values, per model and per
seed, are in \texttt{results/ton\_iot\_dropped\_class\_bound.json}.

Because most of the NF-UNSW-NB15 margin is not architectural, it repays decomposition.
Table~\ref{tab:decomposition} walks from the strongest published configuration to our rungs.
Every level is a three-seed mean, no step carries an interval, and no step is a separation claim.

Figure~\ref{fig:decomposition} draws the walk. Roughly 45\% of the naive gap is our choice of
how to tune, roughly 30\% is our node features,
and roughly 25\% is the architecture. Measured to the fusion model instead of the feedback model, the architecture
share is 15.1\%.

\clearpage
\section{Reproducibility}
\label{app:repro}

This appendix first records the software environment and execution settings, then provides the reproduction commands, and finally identifies the supporting artifacts.

\paragraph{Environment.} Python 3.14.3 in a virtual environment built from
\texttt{requirements.txt}. Message passing, the fold masks and the graph containers come from
PyTorch Geometric \citep{fey2019pyg}. The stratified fold generator and every reported metric
are scikit-learn's \citep{pedregosa2011sklearn}. The ZCA whitener is a local
eigendecomposition rather than a library call. All runs export \texttt{OMP\_NUM\_THREADS=1}. Without that pin,
pooled macro-F1 drifts by roughly 0.012 at a fixed seed, larger than several of the
differences this report measures. One earlier finding in the project was retracted because
of it.

\paragraph{Computational cost.} Everything in this report was produced on one Apple M1 Pro
laptop with 8 cores, using the MPS backend for semantic encoding and CPU for everything else.
Semantic encoding is one measured per-edge inference cost, requiring a frozen forward pass for each edge: at batch size 32 and a 128-token limit it
encodes 171.2 edges per second on NF-UNSW-NB15 and 160.2 on NF-ToN-IoT, so 5.841 and 6.241
milliseconds per edge, and all 656 and 2{,}127 edges in 3.83 and 13.28 seconds. Loading the
encoder costs a further 3.2 seconds once per process. Training is an offline cost paid once:
one five-fold run of the fusion stage at one seed takes 6.4 seconds on NF-UNSW-NB15 and 17.3
on NF-ToN-IoT. These are graphs of a few hundred to a few thousand edges, so the figures say
nothing about how the encoder would scale to per-flow traffic volumes; they establish that
nothing in this study was compute-bound. Measurements are in
\texttt{results/encoding\_cost.json}.

\paragraph{Seeds and folds.} Three training seeds, 42, 1 and 2. The stratified five-fold
partition is generated once at seed 42 and held fixed across all three, so the pooled
out-of-fold edge set is identical between seeds and the seed-matched bootstrap is a paired
comparison. Fold-partition variance is therefore not measured anywhere in this report.

\paragraph{Pipeline.} Graph construction runs from
\texttt{src.pipeline.step1.run\_step1}, writing \texttt{pyg\_data.pt} and
\texttt{aggregated\_edges.csv} under \texttt{data/<dataset>/processed/step1/}. The downstream
stages run through one entry point:

\begin{quote}
\texttt{python -m src.pipeline.common.run\_pipeline --dataset <key> \textbackslash} \\
\texttt{\ \ --only fusion oof prototypes top\_k\_sweep feedback ladder}
\end{quote}

\begin{sloppypar}
with \texttt{<key>} either \texttt{unsw\_nb15} or \texttt{ton\_iot}. That runner passes
\texttt{--use-llm-head} to the sweep and feedback stages, so the feedback model consults the
trained head as the contracts describe. Calling \texttt{train\_feedback} without that flag
reproduces the superseded prototype-consultant ladder instead, at the knobs kept under
\texttt{prototype\_superseded} in each dataset's selected-feedback config. The runner's
\texttt{--seed} is the training seed and reaches every stage that trains; the fold partition
has its own \texttt{--split-seed}, which defaults to 42 and should not be changed if the
results are to stay comparable with the ones reported here. Out-of-fold structural embeddings
must be regenerated with \texttt{step3.build\_oof\_gnn\_embeddings} before the fusion stage
whenever the GNN model changes. The fusion stage does not do it, and will otherwise train on
stale embeddings.

\texttt{reproduce\_ladder.py} rebuilds the NF-UNSW-NB15 ladder end to end and compares it
against the committed contract, printing any rung that differs by more than 0.0005.
\end{sloppypar}

\paragraph{Result contracts.} Every number in this report traces to one of the files in
Table~\ref{tab:artifacts}.

\begin{table}[htbp]
\centering
\caption[Committed artifacts behind the report's numbers]{Committed artifacts behind the report's numbers.}
\label{tab:artifacts}
\footnotesize
\begin{tabular}{@{}p{0.56\textwidth}p{0.40\textwidth}@{}}
\toprule
Artifact & What it carries \\
\midrule
\texttt{results/unsw\_nb15\_current.json} (schema 5) & UNSW ladder, ablations, knobs \\
\texttt{results/ton\_iot\_current.json} (schema 7) & ToN ladder, ablations, knobs \\
\texttt{results/cross\_dataset\_comparison.json} (schema 6) & both datasets side by side \\
\texttt{results/multiseed\_ladder\_v2\_head.json} & three-seed ladder, trained-head consultant \\
\texttt{results/multiseed\_ladder\_v2\_legacy.json} & three-seed ladder, prototype \\
\texttt{results/multiseed\_head\_per\_class.json} & per-class F1, three-seed means \\
\texttt{results/*\_sota\_baselines.json} & re-trained baselines and intervals \\
\texttt{results/*\_oracle\_ceiling\_v2\_trained\_head.json} & the isolation test \\
\texttt{results/knob\_selection\_head/} & the selection curves \\
\texttt{results/consultant\_change\_interval.json} & the interval on the consultant change \\
\texttt{results/consultant\_complementarity.json} & the flagged-set and disagreement figures \\
\texttt{results/ton\_iot\_dropped\_class\_bound.json} & the NF-ToN-IoT scoring asymmetry, bounded \\
\texttt{docs/RESULTS\_ARCHIVE.md} & the full record, including retractions \\
\bottomrule
\end{tabular}
\end{table}

\paragraph{Checks and data access.} The recorded test run passed 174 tests and 359 subtests with \texttt{OMP\_NUM\_THREADS=1}. The build regenerates the tables and runs \texttt{scripts/check\_report\_numbers.py}; its exit status should also be checked directly because the build treats audit failures as warnings. Writing statistics are available separately through \texttt{scripts/prose\_stats.py}. Raw dataset CSVs are not included in the repository. Recalculating the churn diagnostic additionally requires the local traces identified in Appendix~\ref{app:isolation-detail}.
